\section{Introduction}
\label{sec:intro}

Modern deep learning success is built on the assumption of abundant, diverse training data, yet many high-value applications -- rare disease diagnosis, endangered-species monitoring, low-resource robotics -- can only offer a handful of labeled examples per class~\cite{fawakherji2026learning,wang2020generalizing}. Few-shot learning addresses this data bottleneck, and among its solutions, gradient-based meta-learning~\cite{schmidhuber1987evolutionary,bengio2013optimization}, epitomized by Model-Agnostic Meta-Learning (MAML)~\cite{finn2017model}, has become dominant precisely because it is architecture-agnostic: a shared initialization $\theta_0$ is meta-trained so that a handful of gradient steps on a new task's \emph{support set} $S_i$ suffice to adapt it.

That adaptation, however, is a black box. What prior knowledge in $\theta_0$ actually drives the few gradient steps taken on $S_i$, and whether that drive helps or hurts, is not observable from the outside. This opacity is not merely academic: when $S_i$ is hard, noisy, or drawn from a distribution the meta-training tasks never anticipated, adaptation can actively \emph{degrade} query performance -- a failure known as \emph{negative transfer}~\cite{yue2020interventional,si2025meta}. Existing accounts of \emph{why} MAML adapts the way it does are themselves in tension: Raghu~\etal~\cite{raghu2019rapid} argue that most of its benefit comes from reusing near-frozen backbone features and only re-fitting the classifier head (ANIL), whereas Oh~\etal~\cite{oh2020boil} show that under domain shift the opposite regime -- backbone-only adaptation -- dominates (BOIL). Both accounts, and the negative-transfer literature built on them~\cite{yao2019hierarchically,yao2020automated,chen2021hetmaml,parast2026hsfm}, characterize adaptation only through its \emph{symptom}: a change in accuracy or in an intervened representation. None of them explains, at the level of individual support-image regions, \emph{where in feature space} a given adaptation went right or wrong.

Meanwhile, post-hoc explainable AI has matured considerably for static, converged models: Grad-CAM~\cite{selvaraju2020grad}, SHAP~\cite{lundberg2017unified}, and related attribution methods~\cite{simonyan2014visualising,sundararajan2017axiomatic,ribeiro2016lime} all extract a saliency signal from \emph{one} fixed parameter state. This is precisely what a bi-level optimizer like MAML does not have: its meta-test behavior is the endpoint of a $T$-step trajectory $\theta_0 \!\to\! \phi_1 \!\to\! \cdots \!\to\! \phi_T$, and none of these methods aggregate a signal across it. The one method that reasons about meta-training influence at the task level, TL-XML~\cite{mitsuka2025tlxml}, requires access to the full meta-training corpus and still treats the adaptation step itself as an opaque unit -- infeasible whenever meta-training data is proprietary or discarded, and uninformative about \emph{which pixels} of $S_i$ mattered.

This gap motivates two questions we address directly: \emph{(i)} for a given target task, is the adaptation step itself helpful or harmful, and by how much; and \emph{(ii)} which support-set features, located in the shared feature space $F$ that $\theta_0$ and the adapted parameters both live in, are responsible? We answer \emph{(i)} with \textbf{adaptation gain} $\Delta M$, a differentiable, sign-aware scalar contrasting full adaptation against a feature-reuse baseline that freezes the backbone; and \emph{(ii)} with \textbf{\fama} (Feature-space Adjoint Meta-learning Attribution), which back-propagates $\Delta M$ through the \emph{entire} adaptation trajectory using an adjoint recursion, avoiding the $\mathcal{O}(TP^2)$ cost of literal backpropagation-through-time. Because $S_i$ is available at meta-\emph{test} time by construction, \fama needs no access to the meta-training set at all.

In summary, our contributions are:
\begin{itemize}[leftmargin=*,itemsep=1pt,topsep=2pt]
    \item \textbf{Adaptation gain $\Delta M$} (\Cref{sec:adaptation_gain}), a differentiable measure of the benefit or harm caused by letting the backbone adapt on $S_i$, which we prove reduces, to leading order in the fast-adaptation regime, to a term that depends \emph{only} on the backbone's displacement (\Cref{prop:decomposition}).
    \item \textbf{\fama} (\Cref{sec:fama}), a post-hoc explainer that extends Grad-CAM across a multi-step adaptation trajectory via a Pearlmutter Hessian-vector-product adjoint recursion, at $\mathcal{O}(TP)$ cost, requiring no meta-training data and no architectural change.
    \item Two \textbf{trajectory-aware faithfulness protocols} (\Cref{sec:protocols}) -- a bidirectional deletion test and cascading-randomization / support-intervention sanity checks -- that adapt static-model evaluation principles~\cite{samek2016evaluating,petsiuk2018rise,adebayo2018sanity} to a setting where the object of explanation is a trajectory rather than a single prediction; neither protocol is specific to \fama, so both apply unchanged to any future explainer for fast adaptation.
    \item An empirical study on MAML/Conv4 (\Cref{sec:experiments}) that validates \fama quantitatively on tieredImageNet, and, on 600 miniImageNet/tieredImageNet$\to$CUB-200-2011 cross-domain tasks, surfaces \emph{contextual bias} as a concrete, visualizable candidate driver of negative transfer.
\end{itemize}
